% Generated by R/Project_code.R -- auto-generated draft, edit before use
\begin{itemize}
  \item Data: 1575 consecutive days (2013-01-01 to 2017-04-24) after cleaning; 10 invalid pressures interpolated and 33 extreme wind speeds capped at 17.6 km/h.
  \item Strongest relation: meantemp $\sim$ meanpressure, r = -0.881. A straight line explains 77.6\% of the variance of meantemp ($R^2$ = 0.776); the cubic curve raises this to 81.9\% ($R^2$ = 0.819, 2017 test $R^2$ = 0.715).
  \item Strong relations ($R^2$ $\geq$ 0.50): meantemp $\sim$ meanpressure ($R^2$ 0.819, cubic).
  \item Moderate relations ($R^2$ 0.25-0.50): meantemp $\sim$ humidity ($R^2$ 0.331, cubic).
  \item Weak relations ($R^2$ 0.10-0.25): humidity $\sim$ meanpressure ($R^2$ 0.187, quadratic); humidity $\sim$ wind\_speed ($R^2$ 0.165, exponential).
  \item Very weak relations ($R^2$ \textless{} 0.10): meantemp $\sim$ wind\_speed ($R^2$ 0.094, exponential); wind\_speed $\sim$ meanpressure ($R^2$ 0.091, cubic).
  \item Non-linear curves improve meaningfully ($\Delta R^2$ \textgreater{} 0.05) only for: humidity $\sim$ meanpressure: quadratic $\Delta R^2$ = +0.072 (test $R^2$ 0.133 $\rightarrow$ 0.292, confirmed out of sample).
  \item Mild improvement (0.01 \textless{} $\Delta R^2$ $\leq$ 0.05): meantemp $\sim$ meanpressure: cubic $\Delta R^2$ = +0.043 (test $R^2$ 0.755 $\rightarrow$ 0.715, NOT confirmed out of sample); meanpressure $\sim$ meantemp (rev.): quadratic $\Delta R^2$ = +0.012 (test $R^2$ 0.730 $\rightarrow$ 0.752, confirmed out of sample); meanpressure $\sim$ humidity (rev.): cubic $\Delta R^2$ = +0.010 (test $R^2$ -0.213 $\rightarrow$ -0.211, confirmed out of sample).
  \item No meaningful improvement ($\Delta R^2$ $\leq$ 0.01; a straight line is adequate): meantemp $\sim$ humidity; humidity $\sim$ meantemp (rev.); meantemp $\sim$ wind\_speed; wind\_speed $\sim$ meantemp (rev.); humidity $\sim$ wind\_speed; wind\_speed $\sim$ humidity (rev.); wind\_speed $\sim$ meanpressure; meanpressure $\sim$ wind\_speed (rev.).
  \item Possible overfitting: meantemp $\sim$ meanpressure (cubic, $\Delta R^2$ = +0.043, test $R^2$ 0.715 vs 0.755 linear). Here the polynomial gains in-sample but predicts the 2017 test period worse than the straight line, so its higher $R^2$ must be read with caution.
  \item In-sample gains that the selection rule rejected (not supported on the test period): humidity $\sim$ meantemp (rev.) (cubic +0.043); meantemp $\sim$ wind\_speed (cubic +0.013); wind\_speed $\sim$ meantemp (rev.) (cubic +0.012); humidity $\sim$ wind\_speed (cubic +0.011); meanpressure $\sim$ wind\_speed (rev.) (cubic +0.014).
  \item Quadratic and cubic curves bend near the edges of the data, where few days lie (e.g. the cubic temperature-pressure curve peaks at 998.0 hPa and falls again towards the lowest pressures); they describe the observed range only and must not be extrapolated.
  \item The 2017 test period covers only January-April (114 days, one season), so test $R^2$ values are noisy and can be negative even for real relations; they are a check, not a definitive ranking.
  \item Influential points: removing the days with Cook's D \textgreater{} 4/n (42-80 days per model) changes the linear $R^2$ by +0.020 to +0.044. $R^2$ rises in every case: the influential days are poorly fitted days that LOWER the fit, so no reported $R^2$ is inflated by a few outliers.
  \item Wind capping (pre-processing): for meantemp $\sim$ wind\_speed, humidity $\sim$ wind\_speed, wind\_speed $\sim$ meanpressure the linear $R^2$ is 0.098 / 0.160 / 0.090 with capped wind but 0.082 / 0.139 / 0.074 with the raw values (difference +0.016 / +0.021 / +0.015). The capping of 33 extreme values raises these $R^2$ values slightly, so part of the (weak) wind relations depends on that pre-processing decision; the ranking and the conclusion 'weak' do not change.
  \item Seasonality: the calendar month alone explains meantemp 87\%, humidity 58\%, wind\_speed 17\%, meanpressure 88\% of the daily variance ($R^2$ of the monthly means), so temperature and pressure are dominated by the annual cycle.
  \item Share of each linear $R^2$ that disappears when the annual cycle is removed (anomaly correlation): meantemp $\sim$ humidity 8\% (r -0.57 $\rightarrow$ -0.55); meantemp $\sim$ wind\_speed 81\% (r 0.31 $\rightarrow$ 0.14); meantemp $\sim$ meanpressure 76\% (r -0.88 $\rightarrow$ -0.43); humidity $\sim$ wind\_speed 38\% (r -0.40 $\rightarrow$ -0.32); humidity $\sim$ meanpressure 78\% (r 0.34 $\rightarrow$ 0.16); wind\_speed $\sim$ meanpressure 88\% (r -0.30 $\rightarrow$ -0.10).
  \item Temperature-pressure: r = -0.881 on the raw data but -0.431 on the anomalies, so 76\% of the $R^2$ is the shared annual cycle. Within seasons r = -0.51 (Winter), -0.86 (Summer), -0.31 (Monsoon), -0.77 (Post-monsoon). The very strong overall relation is therefore largely seasonal (hot, low-pressure summer/monsoon vs cool, high-pressure winter); a real but weaker day-to-day link remains.
  \item Temperature-humidity: r = -0.574 on the raw data and -0.550 on the anomalies, so only 8\% of the $R^2$ is due to the annual cycle; within seasons r = -0.47 (Winter), -0.81 (Summer), -0.80 (Monsoon), -0.16 (Post-monsoon). Unlike temperature-pressure, this relation is NOT mainly seasonal: the correlation is negative within every season (hotter days are drier), strongly in summer and monsoon. Mixing the seasons actually weakens the overall correlation, because monsoon days are both hot and humid (on average 31.4 degC, 63\% humidity), which puts them on a separate band of the scatter plot; that is why no single curve in humidity reaches a high $R^2$.
  \item Humidity-pressure: r = 0.339 on the raw data but 0.157 on the anomalies (78\% of the $R^2$ is seasonal). The U-shaped quadratic fit reflects the order of the seasons along the pressure axis (humid low-pressure monsoon, dry pre-monsoon, humid high-pressure winter) rather than a physical law, so the curve must not be extrapolated.
  \item Correlation does not imply causation: the regressions describe how the parameters co-vary, not that one causes the other. Temperature, pressure and humidity are all driven by common factors (solar heating, the monsoon circulation), and the choice of response variable is a modelling convention, not a causal direction.
  \item Independence: Durbin-Watson statistics of the linear fits are 0.07-1.13 (2 = no autocorrelation), so the residuals are strongly autocorrelated. The reported p-values and confidence intervals are therefore too optimistic; the $R^2$ values remain valid as descriptive measures of fit.
\end{itemize}
